\documentclass[11pt]{article}
\usepackage[margin=1in]{geometry}
\usepackage{amsmath}
\usepackage{booktabs}
\usepackage{array}
\usepackage{graphicx}
\usepackage{float}
\usepackage{multicol}
\usepackage{xcolor}
\usepackage{hyperref}
\hypersetup{colorlinks=true,linkcolor=blue!45!black,urlcolor=blue!45!black,
  pdftitle={Predicting Glanders Disease Across a Lifestyle or Geographical Area},
  pdfauthor={GlandersWatch Project Team},
  pdfsubject={Chained multi-level machine-learning evaluation}}
\graphicspath{{../figures/}}
\setlength{\parskip}{0.4em}
\sloppy

\title{\textbf{Predicting Glanders Disease Across a Lifestyle or Geographical Area}\\[0.3em]
\large Chained Multi-Level Machine Learning on Four Unused Datasets}
\author{GlandersWatch Project Team}
\date{\today}

\begin{document}
\maketitle

\noindent
\textbf{Summary.} A country-year panel is built from four public sources never previously
joined to glanders: a text-mined Europe PMC occurrence corpus, FAOSTAT livestock stocks, World
Bank indicators and NASA POWER climatology. Six model families are trained under three validation
schemes and evaluated as a three-stage chained pipeline. On a forward temporal holdout the best
model (Random Forest) reaches PR-AUC 0.554 against a
5.6\% base rate, a 10$\times$ lift over chance, with calibration error 0.016.
Chaining does not improve discrimination; climatology transfers poorly across countries; and skill
is weakest in low-income regions.

% =====================================================================

% =====================================================================
\section{Related Work}
\label{sec:related}
% =====================================================================

\paragraph{Glanders epidemiology and burden.}
Glanders is reviewed as a neglected zoonosis of equids with high case fatality and a
persistent risk of human infection \cite{khan2013,kanani2026}. A 2026 systematic review
and meta-analysis consolidates seroprevalence estimates and identifies management and
movement as the dominant epidemiological drivers \cite{purnama2026}. Regionally, Brazil
maintains an active research and control literature \cite{mota2022}, and glanders is
endemic or repeatedly reported across parts of Africa, Asia and the Middle East.

\paragraph{Surveillance in India.}
Indian surveillance is the strongest available evidence base. Malik et al.\ described
outbreaks from 2006 to 2011 \cite{malik2012}; Singha et al.\ reported serological
surveillance of indigenous equines from 2015 to 2018 \cite{singha2020}; and the
2019--2023 national programme screened a large multi-state sample, identifying substantial
geographic concentration of seropositives \cite{singha2026}. These studies establish where
glanders occurs and which districts carry the burden. None of them produces a predictive,
calibrated risk model, and none uses population or reanalysis covariates.

\paragraph{Geographic and spatiotemporal risk analysis.}
Descriptive spatial work exists. Fonseca-Rodr\'iguez et al.\ applied spatiotemporal cluster
detection to Brazilian glanders \cite{fonseca2019}. More generally, species distribution
and ecological niche models have been compared as tools for livestock disease
epidemics \cite{hollings2017}. This literature is \emph{descriptive}: it identifies
clusters post hoc rather than forecasting future cells.

\paragraph{Machine learning for animal-disease risk prediction.}
Predictive modelling in veterinary epidemiology is established for avian influenza.
Kane et al.\ compared ARIMA and random forest for H5N1 outbreak prediction
\cite{kane2014}; Walsh et al.\ used machine learning to target active surveillance
\cite{walsh2019}; Yoo et al.\ built dynamic farm-level risk prediction in the Republic of
Korea \cite{yoo2022}; and Alshatti et al.\ produced a risk prediction model for HPAI in
Kuwait \cite{alshatti2026}. This is the closest methodological template for the present
work, and it exists for influenza rather than for glanders.

\paragraph{Machine learning and glanders: the gap.}
The only machine-learning work specific to \emph{Burkholderia mallei} that we located is
spectroscopic identification of the organism --- Raman assays classified by learned
models \cite{moawad2019,stoeckel2015}. That is a laboratory diagnostic problem, not a
population-level forecasting problem. A prior-art audit we ran over Europe PMC
(Section~\ref{sec:novelty}) returned zero records combining glanders with NASA POWER
climatology and four combining glanders with gradient boosting, none of which is a risk
model. The gap addressed here is therefore specific: a calibrated, chained,
population-level occurrence model for glanders.

\paragraph{Text mining for disease event extraction.}
Using literature as a structured occurrence signal has precedent in epidemic intelligence
\cite{collier2012}, including multilingual event extraction for outbreak
detection \cite{lejeune2015}. Our contribution is to apply sentence-level occurrence
extraction to build the \emph{label} itself when official event data are incomplete, and
to report the surveillance bias that this induces.

\paragraph{Methodological foundations.}
The model families are the standard tabular ensembles: random forests \cite{breiman2001},
XGBoost \cite{chen2016} and LightGBM \cite{ke2017}. Feature attribution uses
SHAP \cite{lundberg2017}. Because the outcome is rare, precision-recall analysis is the
appropriate primary metric \cite{saito2015}, and the relationship between PR and ROC
curves is well characterised \cite{davis2006}. Probability calibration follows
Niculescu-Mizil and Caruana \cite{niculescu2005}, with the measurement framework of Huang
et al.\ \cite{huang2020}. Validation follows Roberts et al., who show that ignoring
temporal and spatial dependence inflates performance estimates \cite{roberts2017} --- the
reason our primary scheme is a forward temporal holdout and our selection criterion is
leave-country-out.

\paragraph{Data sources.}
Official animal-health reporting is drawn from WOAH WAHIS \cite{wahis}; Indian policy and
case definitions from the 2025 National Action Plan \cite{dahd2025}; livestock
populations from FAOSTAT \cite{faostat}; socio-economic covariates from the World Bank
\cite{wdi}; climatology from NASA POWER \cite{power}; and the literature corpus from
Europe PMC \cite{europepmc}.

\paragraph{Positioning.}
The combination above is, to our knowledge, unstudied. Prior glanders work is either
epidemiological description or laboratory classification; prior veterinary ML risk work is
built for influenza; and prior text-mined occurrence work is not used for glanders at all.
This study joins an occurrence corpus, a livestock census, development indicators and
reanalysis climatology into a single country-year panel, and reports both what predicts and
what does not.


\section{Dataset Used}

\begin{table}[H]
\centering
\caption{Sources. None has previously been used for glanders risk modelling.}
\small
\begin{tabular}{@{}>{\raggedright\arraybackslash}p{6.4cm}>{\raggedright\arraybackslash}p{4.0cm}>{\raggedright\arraybackslash}p{4.4cm}@{}}
\toprule
\textbf{Source} & \textbf{Role} & \textbf{Volume} \\
\midrule
Europe PMC, text-mined occurrences & Animal history & 6,674 records / 263 events \\
FAOSTAT livestock stocks (7 species) & Animal population & 68,449 stock rows, 1961--2023 \\
World Bank WDI + centroids & Socio-economic context & 10 indicators, 217 countries \\
NASA POWER agro-climatology & Environment & 6 variables, 211 centroids \\
\bottomrule
\end{tabular}
\end{table}

\noindent
All sources are keyed on ISO-3 country and calendar year. FAOSTAT is streamed from a 32\,MB zip
($\approx$545\,MB normalised) and filtered to seven stock series. Europe PMC is paged with a cursor
and reduced to sentence-level events by a country-gazetteer plus occurrence-lexicon extractor.
\label{sec:novelty}\emph{Novelty:} a prior-art audit of Europe PMC returns
\textbf{0} records for glanders with NASA POWER and
4 for glanders with gradient boosting, none of which
is a risk model.

\paragraph{Panel.} 6,700 country-year rows, 268
countries, 2000--2024. Modelling is restricted to equine-holding countries (maximum stock
$\geq$1{,}000), years 2005--2023, with an observable next-year label:
\textbf{3,078 rows $\times$ 44 features}, \textbf{172 positives}.

\paragraph{Preprocessing.} Country and area names reconciled to ISO-3 via \texttt{pycountry}
with explicit overrides; World Bank aggregates dropped by requiring a 3-character code; literature
events de-duplicated on \texttt{(iso3, year, record\_id)} so one paper cannot inflate a cell.
Median imputation with explicit missing indicators, fitted inside the pipeline. A country-year is
coded 0 only after the country passes the equine-stock filter, so absence is not inferred from
silence alone.

\paragraph{Feature types and missingness.} Continuous numeric (stocks, climate, economics),
categorical (World Bank region and income), and count/derived history. Missingness is low and
source-specific: FAOSTAT 0.00\%, World Bank
1.68\%, NASA POWER
2.19\%, overall 0.99\%. No column
exceeds 5.8\% missing.

% =====================================================================
\section{Dataset Analysis}

\begin{table}[H]
\centering
\caption{Label anatomy.}
\small
\begin{tabular}{@{}lr@{}}
\toprule
\textbf{Property} & \textbf{Value} \\
\midrule
Observations & 3,078 \\
Positives / negatives & 172 / 2,906 \\
Prevalence & 5.59\% \\
Imbalance ratio & 16.9 : 1 \\
Countries with any positive & 56 / 162 \\
Repeat-positive countries & 30 (54\%) \\
Positives in top-5 countries & 35.5\% \\
\bottomrule
\end{tabular}
\end{table}

\noindent
Three consequences. The 17:1 imbalance rules out accuracy as a metric.
54\% of positive countries fire in $\geq$2 years, so the label is
temporally persistent and history features are legitimately exploitable. Positives concentrate in
top countries, and the highest-risk cells are India, the USA and Australia --- reference-laboratory
capacity, not incidence. The target is therefore treated throughout as
\emph{surveillance-observable activity}, not true incidence.

\paragraph{Association and redundancy.} Spearman $\rho$ and mutual information were computed for
every numeric feature; Cram\'er's $V$ for the categorical strata. The strongest monotone signal is
\texttt{years\_since\_last} ($\rho=-0.35$): recent activity predicts future activity. Variance
inflation reaches 228 for \texttt{temp\_mean},
because annual country-level climate variables are near-collinear; tree ensembles are invariant to
this, so it affects interpretation only.

\paragraph{The selection finding.} Marginal and conditional importance disagree almost completely:
rank correlation between mutual information and SHAP is \textbf{-0.03}, with only
5/10 top features shared. Longitude ranks first marginally but $\approx$20th
conditionally, because once population and history are conditioned on it carries little extra
information. Features were therefore selected on conditional evidence (ablation and SHAP), with
univariate statistics used only to screen zero-signal columns.

\begin{table}[H]
\centering
\caption{Top 8 features by mean $|$SHAP$|$, with marginal statistics.}
\small
\begin{tabular}{@{}rlrrrr@{}}
\toprule
\textbf{\#} & \textbf{Feature} & \textbf{mean $|$SHAP$|$} & \textbf{Spearman $\rho$} &
\textbf{MI} & \textbf{VIF} \\
\midrule
1 & pop\_total & 1.144 & +0.28 & 0.0639 & 4 \\
2 & internet\_pct & 0.542 & +0.09 & 0.0015 & 5 \\
3 & temp\_mean\_anom & 0.469 & -0.00 & 0.0012 & 2 \\
4 & log\_livestock & 0.458 & +0.25 & 0.0571 & 3 \\
5 & humidity\_anom & 0.370 & +0.02 & 0.0066 & 2 \\
6 & years\_since\_last & 0.353 & -0.35 & 0.0486 & 2 \\
7 & wind & 0.346 & -0.08 & 0.0210 & 3 \\
8 & sheep & 0.344 & +0.20 & 0.0512 & 13 \\
\bottomrule
\end{tabular}
\end{table}

\paragraph{Feature creation.} Four groups, all using information at $t$ or earlier.
\emph{History} (7): lags, a 5-year rolling count, years since last event, and world/regional
contemporaneous shares. \emph{Population} (14): log equid and livestock stocks, equid share,
density, 1- and 3-year growth, seven species counts; livestock units use metabolic weights
(cattle and camels 0.8, small ruminants 0.1). \emph{Climate} (10): annual temperature extremes,
precipitation total and coefficient of variation, humidity, wind, plus expanding-mean
\emph{anomalies} --- the anomaly distinguishes ``unusual for this country'' from latitude.
\emph{Socio-geography} (11): income, agricultural intensity, GDP, rural share, internet
penetration, population, density, absolute latitude and longitude.

\paragraph{What was dropped.} Four columns with mutual information below 0.002 and
$|\rho|<0.02$. The raw \texttt{equids} sum, which is exactly \texttt{horses + asses + mules}.
Three redundant clusters at $|\rho|\geq0.90$ were collapsed to one representative each. Scaling is
applied only to the logistic baseline; tree models receive raw values. No PCA: with 44 features the
binding constraint is signal, not dimensionality.

% =====================================================================
\section{Models Used}

\begin{table}[H]
\centering
\caption{Model families and selected hyperparameters.}
\small
\begin{tabular}{@{}>{\raggedright\arraybackslash}p{3.4cm}>{\raggedright\arraybackslash}p{3.1cm}>{\raggedright\arraybackslash}p{8.2cm}@{}}
\toprule
\textbf{Model} & \textbf{Family} & \textbf{Selected configuration} \\
\midrule
Logistic regression & linear & $C=1.0$, $\ell_2$, balanced \\
Decision tree & tree & depth 3, leaf 1 \\
Random forest & bagged trees & depth 6, leaf 1, mtry sqrt \\
Extra trees & bagged trees & depth 6, leaf 10, mtry 0.3 \\
XGBoost & boosted trees & depth 3, $\lambda$=1.0 \\
LightGBM & boosted trees & leaves 15, min child 40 \\
\bottomrule
\end{tabular}
\end{table}

\noindent
The choice is dictated by three properties: a 17:1 imbalance,
172 positives, and collinear continuous features. Neural networks are excluded on
evidence --- with 172 positives and 44 columns no parameter-to-sample ratio is
defensible. Ensembles handle collinearity and non-monotone interactions without manual
specification. The linear model is kept as a floor: if boosting cannot beat it, the nonlinearity is
not earning its complexity.

\paragraph{Selection.} Models were selected on \emph{leave-country-out} PR-AUC, not temporal
holdout, because geographic generalisation is the deployment claim. \textbf{Extra Trees} was
selected (0.363). No stacking ensemble was adopted: the between-family spread is small
relative to the bootstrap interval, and a meta-learner fitted on 172 positives would
add variance without a defensible gain. Probabilities are calibrated instead, which is what actually
matters for threshold-based alerting.

% =====================================================================
\section{Evaluation and Training Pipeline}

\paragraph{Splits.} Three schemes, each answering a different question.
\emph{Temporal holdout} (train 2005--2017, test 2018--2023) asks whether the system can be used
next year, and is the primary scheme. \emph{Leave-country-out} (5 folds grouped by ISO-3) asks
whether it transfers to an unseen country, and drives selection. \emph{Repeated stratified 5-fold}
(3 repeats) is a stability check. Grouping is essential: a random row split would place the same
country in train and test and leak the label through geography.

\paragraph{Metrics.} Precision-recall AUC (average precision) is primary, because at
5.6\% prevalence ROC-AUC is dominated by true negatives and flatters every model.
ROC-AUC is reported for comparability. Brier score and expected calibration error are reported
because an uncalibrated probability cannot support an alert threshold. Precision, recall,
specificity, F1 and MCC use a threshold tuned on \emph{training} data only. Bootstrap percentile
intervals (1{,}000 resamples) quantify sampling uncertainty.

\paragraph{Leakage controls.} Features are built from $t$ and earlier; only \texttt{y\_*}
columns use $t+1$. Splits are chronological. Hyperparameters are tuned inside training folds via
\texttt{GroupKFold} on country. Calibration is fitted on training folds only. The alert threshold
is fixed on the training period and applied unchanged to test. Missing indicators are generated
inside the pipeline, so the test set never influences imputation.

\paragraph{Workflow and tuning.} Ingest (resumable, cached) $\rightarrow$ assemble panel with
lags, rolling windows, anomalies and shifted targets $\rightarrow$ analyse $\rightarrow$
\texttt{RandomizedSearchCV} with 8 draws scored on average precision, grouped by country,
refit on the full fold $\rightarrow$ calibrate $\rightarrow$ evaluate under all three schemes
$\rightarrow$ chain the three stages. The search space is deliberately narrow (40--60 candidates
per family) because with 172 positives a wide search overfits the validation
objective.

\paragraph{Compute and reproducibility.} Full run: 4.5 minutes,
CPU only, single process, Apple M4 (10 cores, 16\,GB). Per-family fit times range from
1.9\,s (logistic) to
50.7\,s (LightGBM). One seed constant (42) governs
fitting, CV, bootstrap and SHAP sampling. Raw payloads are cached and the panel is versioned on
disk, so every number regenerates without re-downloading. Caveat: \texttt{n\_jobs=-1} means
ensemble scores may differ in the last decimal across thread counts. Environment: Python
3.12.13, scikit-learn 1.9.1, XGBoost 3.4.1,
LightGBM 4.7.0, SHAP 0.54.1.

% =====================================================================
\section{Results}

\begin{table}[H]
\centering
\caption{Level 1 (country-year). Brackets are 95\% bootstrap intervals on the temporal holdout.
$\star$ marks the model selected on leave-country-out PR-AUC.}
\small
\resizebox{\textwidth}{!}{%
\begin{tabular}{@{}lccccccc@{}}
\toprule
\textbf{Model} & \textbf{PR-AUC [CI]} & \textbf{ROC-AUC} & \textbf{PR-AUC LOCO} &
\textbf{PR-AUC CV} & \textbf{Brier} & \textbf{Prec} & \textbf{Rec} \\
\midrule
Logistic regression & 0.454 [0.34,0.58] & 0.812 & 0.223 & 0.414 & 0.0547 & 0.25 & 0.57 \\
Decision tree & 0.389 [0.29,0.50] & 0.861 & 0.267 & 0.360 & 0.0602 & 0.49 & 0.53 \\
Random forest & 0.554 [0.43,0.66] & 0.870 & 0.304 & 0.476 & 0.0503 & 0.53 & 0.50 \\
Extra trees$^{\star}$ & 0.530 [0.42,0.64] & 0.854 & 0.363 & 0.476 & 0.0535 & 0.47 & 0.51 \\
XGBoost & 0.485 [0.37,0.60] & 0.844 & 0.281 & 0.423 & 0.0536 & 0.61 & 0.38 \\
LightGBM & 0.496 [0.38,0.61] & 0.841 & 0.339 & 0.397 & 0.0534 & 0.65 & 0.30 \\
\bottomrule
\end{tabular}}
\end{table}

\noindent
At 5.6\% prevalence the best temporal PR-AUC is a \textbf{10$\times$ lift} over
chance. Every nonlinear ensemble beats the linear baseline and the single tree is worst, so the
relationships are non-additive. The more important pattern is the drop under leave-country-out:
roughly 40\% relative for the ensembles, which is the honest estimate of transfer to a new country.
Regression degrades least in absolute terms, suggesting a meaningful share of the transferable
signal is additive and population-driven.

\begin{figure}[H]
\centering
\includegraphics[width=0.95\textwidth]{fig02_model_comparison.png}
\caption{Discrimination by model and scheme. The gap between repeated CV and leave-country-out is
the cost of generalising to an unseen country.}
\end{figure}

\paragraph{Calibration.} Expected calibration error \textbf{0.016},
Brier 0.0535. A stated 24\% risk corresponds to roughly 24\% observed
frequency, which is what makes threshold-based alerting defensible.

\paragraph{Coarse resolution.} Aggregating to country $\times$ 5-year periods raises prevalence
to 14.7\% and improves every model, isolating the mechanism: coarser windows match the
label's own reporting cadence.

\begin{table}[H]
\centering
\caption{Ablation settings (XGBoost, fixed parameters).}
\small
\begin{tabular}{@{}lccc@{}}
\toprule
\textbf{Setting} & \textbf{$k$} & \textbf{PR-AUC (time)} & \textbf{PR-AUC (LOCO)} \\
\midrule
AS-1 history only & 7 & 0.407 & 0.174 \\
AS-2 population only & 14 & 0.432 & 0.241 \\
AS-3 climate only & 10 & 0.395 & 0.078 \\
AS-5 history+population & 21 & 0.451 & 0.321 \\
AS-6 full (all groups) & 42 & 0.487 & 0.259 \\
\bottomrule
\end{tabular}
\end{table}

\noindent
Two findings matter more than the headline. \emph{Parsimony generalises better}: the 21-feature
history-plus-population set scores 0.321 LOCO
against 0.259 for the full 42-feature set, while
scoring \emph{lower} temporally. \emph{Climate alone barely transfers}:
0.078 against a 0.056 floor --- climatology at
annual country resolution mostly encodes latitude, which is country identity in disguise.

\paragraph{The chain.} Stage 1 forecast next-year log equid stock (MAE 0.424,
Spearman 0.976); stage 3 predicted log event counts among positives (MAE
0.150, Spearman 0.465); the hurdle
product ranked annual burden at Spearman 0.326.
\textbf{Chaining added nothing to discrimination}: -0.003 PR-AUC, inside noise.
The reason is structural --- population level is already a direct feature, so forecasting it and
feeding the forecast back is redundant. The chain earns its place at the alert layer instead, where
its top decile captures 70 of 106
events.

\begin{figure}[H]
\centering
\includegraphics[width=0.95\textwidth]{fig07_alert_budget.png}
\caption{Alert economics. Roughly 6\% of country-years must be alerted to reach 50\% recall;
28\% reaches 80\%, against 6\% for random alerting.}
\end{figure}

\noindent
This is the operationally meaningful result: alerting 6\% of country-years captures half of future
documented events, roughly eight times the efficiency of arbitrary selection.

\paragraph{Where it fails.}

\begin{table}[H]
\centering
\caption{Leave-country-out performance by region.}
\small
\begin{tabular}{@{}lrrcc@{}}
\toprule
\textbf{Region} & \textbf{$n$} & \textbf{events} & \textbf{ROC-AUC} & \textbf{PR-AUC} \\
\midrule
East Asia \& Pacific & 456 & 48 & 0.865 & 0.396 \\
Europe \& Central Asia & 912 & 31 & 0.826 & 0.163 \\
Latin America \& Caribbean & 551 & 12 & 0.903 & 0.381 \\
MENA, Afghanistan \textbackslash{}\& Pakistan & 380 & 23 & 0.834 & 0.220 \\
Sub-Saharan Africa & 665 & 27 & 0.699 & 0.073 \\
\bottomrule
\end{tabular}
\end{table}

\noindent
Skill is lowest in Sub-Saharan Africa and in low-income countries --- precisely where surveillance
is thinnest and an early-warning system would matter most. The mechanism is label sparsity, not
model failure. These strata should not be treated as interchangeable in a deployment decision.

\paragraph{Caveats.} The label is documented activity, and a zero conflates absence of disease
with absence of publication; all scores describe skill at predicting surveillance-observable events.
Text-mined labels inherit extraction error and the extractor has not been validated against a
hand-annotated gold set. Country-level aggregation hides the district heterogeneity that matters in
India. With 172 positives, intervals are wide ($\approx\pm$0.11 on the headline
PR-AUC) and subgroup estimates are noisy. Geographic and socio-economic features act largely as
reporting-capacity proxies, so no causal or intervention claim follows.

\section*{Reproducibility}
\addcontentsline{toc}{section}{Reproducibility}

\noindent
Scripts \texttt{src/00} -- \texttt{src/08} run in order. Cached payloads in
\texttt{data/raw}; assembled panel and events in \texttt{data/processed}; all reported numbers
in \texttt{reports/results.json} and \texttt{reports/dataset\_analysis.json}. Every table and
figure here is generated from those files --- the report contains no hand-transcribed results.

\begin{multicols}{2}
\begin{thebibliography}{99}
\footnotesize
\setlength{\itemsep}{0.15em}
\setlength{\parskip}{0pt}
\bibitem{khan2013} Khan I, Wieler LH, Melzer F, Elschner MC, Muhammad G, Ali S, Sprague LD, Neubauer H, Saqib M. Glanders in animals: a review on epidemiology, clinical presentation, diagnosis and countermeasures. \emph{Transboundary and Emerging Diseases}. 2013;60(3):204--221. doi:10.1111/j.1865-1682.2012.01342.x
\bibitem{kanani2026} Kanani Y, Oryan A, Derakhshandeh A. Glanders: past, present and future of a neglected zoonosis as a threat in veterinary and human health. \emph{Veterinary Medicine and Science}. 2026. doi:10.1002/vms3.70950
\bibitem{purnama2026} Purnama MTE, Firdausy LW, Fikri F, Wicaksono AP, Maslamama ST. Seroprevalence and associated epidemiological drivers of glanders in horses: a systematic review and meta-analysis. \emph{Research in Veterinary Science}. 2026;106416. doi:10.1016/j.rvsc.2026.106416
\bibitem{mota2022} Mota RA, Pinheiro J\'unior JW. Current status of glanders in Brazil: recent advances and challenges. \emph{Brazilian Journal of Microbiology}. 2022. doi:10.1007/s42770-022-00814-1
\bibitem{singha2026} Singha H, Shanmugasundaram K, Jhandai P, Rani R, Yashpal, Gupta P, et al. Surveillance of equine glanders in India (2019--2023): prevalence estimate, risk analysis and future strategy. \emph{Microbial Pathogenesis}. 2026;217:108597. doi:10.1016/j.micpath.2026.108597
\bibitem{singha2020} Singha H, Shanmugasundaram K, Tripathi BN, Saini S, Khurana SK, Kanani A, et al. Serological surveillance and clinical investigation of glanders among indigenous equines in India from 2015 to 2018. \emph{Transboundary and Emerging Diseases}. 2020. doi:10.1111/tbed.13475
\bibitem{malik2012} Malik P, Singha H, Khurana SK, Kumar R, Kumar S, Raut AA, et al. Emergence and re-emergence of glanders in India: a description of outbreaks from 2006 to 2011. \emph{Veterinaria Italiana}. 2012;48(2):167--178. PMID:22718333
\bibitem{fonseca2019} Fonseca-Rodr\'iguez O, Pinheiro J\'unior JW, Mota RA. Spatiotemporal analysis of glanders in Brazil. \emph{Journal of Equine Veterinary Science}. 2019. doi:10.1016/j.jevs.2019.03.216
\bibitem{hollings2017} Hollings T, Robinson A, van Andel M, Jewell C, Burgman M. Species distribution models: a comparison of statistical approaches for livestock and disease epidemics. \emph{PLOS ONE}. 2017;12(8):e0183626. doi:10.1371/journal.pone.0183626
\bibitem{alshatti2026} Alshatti M, Al-Hemoud A, Othman A, Al-Khalaifah H. A risk prediction model for highly pathogenic avian influenza outbreaks in Kuwait. \emph{Frontiers in Veterinary Science}. 2026. doi:10.3389/fvets.2026.1843962
\bibitem{kane2014} Kane MJ, Price N, Scotch M, Rabinowitz P. Comparison of ARIMA and Random Forest time series models for prediction of avian influenza H5N1 outbreaks. \emph{BMC Bioinformatics}. 2014;15:276. doi:10.1186/1471-2105-15-276
\bibitem{walsh2019} Walsh DP, Ma TF, Ip HS, Zhu J. Artificial intelligence and avian influenza: using machine learning to enhance active surveillance for avian influenza viruses. \emph{Transboundary and Emerging Diseases}. 2019. doi:10.1111/tbed.13318
\bibitem{yoo2022} Yoo DS, Song YH, Choi DW, Lim JS, Lee K, Kang T. Machine learning-driven dynamic risk prediction for highly pathogenic avian influenza at poultry farms in the Republic of Korea. \emph{Transboundary and Emerging Diseases}. 2021. doi:10.1111/tbed.14419
\bibitem{moawad2019} Moawad AA, Silge A, Bocklitz T, Fischer K, R\"osch P, Roesler U, Elschner MC, Popp J, Neubauer H. A machine learning-based Raman spectroscopic assay for the identification of \emph{Burkholderia mallei} and related species. \emph{Molecules}. 2019;24(24):4516. doi:10.3390/molecules24244516
\bibitem{stoeckel2015} St\"ockel S, Meisel S, Elschner M, Melzer F, R\"osch P, Popp J. Raman spectroscopic detection and identification of \emph{Burkholderia mallei} and \emph{Burkholderia pseudomallei}. \emph{Analytical and Bioanalytical Chemistry}. 2015;407(3):833--843. doi:10.1007/s00216-014-7906-5
\bibitem{collier2012} Collier N. Uncovering text mining: a survey of current work on web-based epidemic intelligence. \emph{Global Public Health}. 2012;7(7):731--749. doi:10.1080/17441692.2012.699975
\bibitem{lejeune2015} Lejeune G, Brixtel R, Doucet A, Lucas N. Multilingual event extraction for epidemic detection. \emph{Artificial Intelligence in Medicine}. 2015;65(2):131--143. doi:10.1016/j.artmed.2015.06.005
\bibitem{chen2016} Chen T, Guestrin C. XGBoost: a scalable tree boosting system. In: \emph{Proceedings of the 22nd ACM SIGKDD International Conference on Knowledge Discovery and Data Mining}. 2016:785--794. doi:10.1145/2939672.2939785
\bibitem{ke2017} Ke G, Meng Q, Finley T, Wang T, Chen W, Ma W, Ye Q, Liu TY. LightGBM: a highly efficient gradient boosting decision tree. In: \emph{Advances in Neural Information Processing Systems 30 (NIPS)}. 2017:3146--3154.
\bibitem{breiman2001} Breiman L. Random forests. \emph{Machine Learning}. 2001;45(1):5--32. doi:10.1023/A:1010933404324
\bibitem{lundberg2017} Lundberg SM, Lee SI. A unified approach to interpreting model predictions. In: \emph{Advances in Neural Information Processing Systems 30 (NIPS)}. 2017:4765--4774. arXiv:1705.07874
\bibitem{saito2015} Saito T, Rehmsmeier M. The precision-recall plot is more informative than the ROC plot when evaluating binary classifiers on imbalanced datasets. \emph{PLOS ONE}. 2015;10(3):e0118432. doi:10.1371/journal.pone.0118432
\bibitem{davis2006} Davis J, Goadrich M. The relationship between precision-recall and ROC curves. In: \emph{Proceedings of the 23rd International Conference on Machine Learning}. 2006:233--240. doi:10.1145/1143844.1143874
\bibitem{roberts2017} Roberts DR, Bahn V, Ciuti S, Boyce MS, Elith J, Guillera-Arroita G, et al. Cross-validation strategies for data with temporal, spatial, hierarchical, or phylogenetic structure. \emph{Ecography}. 2017;40(8):913--929. doi:10.1111/ecog.02881
\bibitem{niculescu2005} Niculescu-Mizil A, Caruana R. Predicting good probabilities with supervised learning. In: \emph{Proceedings of the 22nd International Conference on Machine Learning}. 2005:625--632. doi:10.1145/1102351.1102430
\bibitem{huang2020} Huang Y, Li W, Macheret F, Gabriel RA, Ohno-Machado L. A tutorial on calibration measurements and calibration models for clinical prediction models. \emph{Journal of the American Medical Informatics Association}. 2020;27(4):621--633. doi:10.1093/jamia/ocz228
\bibitem{wahis} World Organisation for Animal Health (WOAH). World Animal Health Information System (WAHIS). \url{https://wahis.woah.org} (accessed November 2026).
\bibitem{dahd2025} Department of Animal Husbandry and Dairying, Government of India. \emph{National Action Plan for Control and Eradication of Glanders in India} (Revised 2025). New Delhi: DAHD; 2025.
\bibitem{faostat} Food and Agriculture Organization of the United Nations. FAOSTAT: production -- crops and livestock products. \url{https://www.fao.org/faostat} (accessed November 2026).
\bibitem{power} NASA Langley Research Center. Prediction Of Worldwide Energy Resources (POWER) project: agroclimatology API. \url{https://power.larc.nasa.gov} (accessed November 2026).
\bibitem{wdi} World Bank. World Development Indicators. \url{https://databank.worldbank.org} (accessed November 2026).
\bibitem{europepmc} Europe PMC. Literature search API. \url{https://europepmc.org} (accessed November 2026).
\end{thebibliography}
\end{multicols}

\end{document}
